\section{敏感性分析}

\subsection{随机场景稳健性（seed 波动）}

同一布局在 seed $0$--$19$ 的 20 局中虚拟时间 $5580$--$7861$ s（均值 $6445$ s、极差
$2281$ s），清除率 $256/256$ 全部完成、误差全部 $\le 20$ m——策略对场景随机性稳健，
结果按 seed 逐字节可复现。

\subsection{接收半径 $R$ 的对抗}

$R\in[1000,1500]$ m 连续未知，是本文对"漏听"最敏感的物理量：\textbf{贴边对抗}（源贴
$1770$ m、波束径向 $\pm 8^\circ$、$R$ 取三档 $1000/1250/1500$）$43\,200$ 例\textbf{0 漏}
，说明布局对 $R$ 全区间稳健；残余漏例全部出现在"内部半径源 $+$ 朝外波束 $+$ 方向恰落
7 基点方位空隙"（精细对抗 $181\,440$ 例漏 $107$），与 $R$ 无单调关系。

\subsection{布局参数的灵敏度}

\begin{itemize}
  \item \textbf{外圈半径}：$1850$ m（贴边源迎光区内沿 $1770$ m 之内的近邻）vs 中点外推
        $2270$ m——后者离内部源迎光面太远，听率仅 $99.81\%$；外圈比 $1850$ m 再缩，则
        覆盖不到贴边源迎光区外沿。$1850$ m 是贴边与内部两类源迎光面覆盖的平衡点。
  \item \textbf{外圈个数}：$12$（间隔 $30^\circ$）vs $11$（$32.7^\circ$）——$11$ 点听率
        略降（$99.989\% \to 99.986\%$）但顺路清除扇区变宽、20 局总里程反升 $464$ m；$12$
        为实战极值。
  \item \textbf{内部补点开关}：加 $3$ 内点（23 点版）听率 $99.9974\%$ 但 20 局慢 $243$ s；
        20 点版听率 $99.9891\%$ 且 $256/256$ 全清——按"任务完成率 + 耗时"双标准，20 点
        为稳健最优。
  \item \textbf{NN 精英精评的采样墙}：$10$ 万级蒙特卡洛只能分辨 $\approx 0.03\%$ 级的听率
        差异，无法逼近人工 $400$ 万例全口径的 $0.01\%$ 级精调——这是采样统计墙（方法固有
        局限），非布局缺陷。
\end{itemize}

\subsection{环境常数}

移动速度、测向/切换/清除耗时、示向度误差均为题目给定或官方引擎固定值；本地模拟器
\texttt{jammers-py} 与官方引擎同判据、同时间常量，离线对拍结论可直接迁移到官方模式
（\texttt{T4.py} 不加参数即连官方模拟器）。